% GENERATED FILE. Edit canonical lessons, then run npm run generate:books.
% canonical-source-hash: fnv1a64:7989e67fef3817a5
% canonical-lessons: RU-C279-izvestiye, RU-C279-vypusk, RU-C279-zagolovok, RU-C279-statya, RU-C279-sensatsiya

\chapter{News and headlines}
\label{ch:ru-news-and-headlines}
\hlchaptermodality{279}

\begin{chapteropening}
\textbf{By the end of this chapter:} \emph{I can talk about the news: a piece of news, a bulletin, a headline, an article and a sensation.}

Complete the last lesson of chapter 279: I can talk about the news: a piece of news, a bulletin, a headline, an article and a sensation.
\end{chapteropening}

\section[\texorpdfstring{\ru{известие} (izvéstiye)}{известие (izvéstiye)}]{\ru{известие} (izvéstiye) — a piece of news, a report}

\label{lesson:RU-C279-izvestiye}



\begin{quote}
Before the new one: say the Russian for to recall, then the Russian for to memorize.
\end{quote}

\subsection*{You'll want to know: \ru{известие}}
\textbf{\ru{известие}} (\emph{izvéstiye}) — \textquotedblleft{}a piece of news, a report\textquotedblright{}.

\begin{quote}
\ru{Мы} \ru{получили} \ru{известие} \ru{о} \ru{его} \ru{приезде}.
\end{quote}

\begin{quote}
\emph{My poluchíli izvéstiye o yevó priyézde.}
\end{quote}

\begin{quote}
We received news of his arrival.
\end{quote}

More formal than \emph{\ru{новость}}. Neuter; plural \emph{\ru{известия}}.

\begin{cousinweb}
\emph{\ru{известие}} is built on \emph{\ru{весть}} \textquotedblleft{}news, tidings\textquotedblright{}, from the old Slavic root of \emph{\ru{ведать}} \textquotedblleft{}to know\textquotedblright{}, which goes back to PIE \emph{weyd-} \textquotedblleft{}to see, to know\textquotedblright{} (EtymDB lists \emph{\ru{известие}} \textquotedblleft{}news\textquotedblright{} with no parent).
\end{cousinweb}

\begin{grammarlens}[title={What you've built}]
One more word for telling, explaining and describing at B1.
\end{grammarlens}

\subsection*{Your turn}
\begin{itemize}
\raggedright
  \item \emph{Say it:} \emph{izvéstiye}
  \item \emph{Say it:} \emph{izvéstiye}, once more
  \item \emph{Say it:} a sentence of your own with \emph{izvéstiye}
  \item \emph{Recall:} say the Russian for to recall, then the Russian for to memorize, then say \emph{izvéstiye} again
\end{itemize}

\begin{tcolorbox}[breakable,colback=teal!4,colframe=teal!35!black,title={Before you move on}]
What does \ru{известие} mean? (\textquotedblleft{}A piece of news, a report\textquotedblright{}.) Say it once more.
\end{tcolorbox}
\section[\texorpdfstring{\ru{выпуск} (výpusk)}{выпуск (výpusk)}]{\ru{выпуск} (výpusk) — a bulletin, an issue, an edition}

\label{lesson:RU-C279-vypusk}



\begin{quote}
Before the new one: say the Russian for to memorize, then the Russian for a piece of news.
\end{quote}

\subsection*{You'll want to know: \ru{выпуск}}
\textbf{\ru{выпуск}} (\emph{výpusk}) — \textquotedblleft{}a bulletin, an issue, an edition\textquotedblright{}.

\begin{quote}
\ru{В} \ru{вечернем} \ru{выпуске} \ru{новостей} \ru{говорили} \ru{о} \ru{погоде}.
\end{quote}

\begin{quote}
\emph{V vechérnem výpuske novostéy govoríli o pogóde.}
\end{quote}

\begin{quote}
In the evening news bulletin they talked about the weather.
\end{quote}

\emph{\ru{Выпуск} \ru{новостей}} — a news bulletin; also an issue of a magazine.

\begin{cousinweb}
\emph{\ru{выпуск}} is built on \emph{\ru{выпускать}} \textquotedblleft{}to let out, to issue\textquotedblright{}, from \emph{\ru{пускать}} \textquotedblleft{}to let go\textquotedblright{}, inherited from Proto-Slavic \emph{pustiti} (EtymDB): what is let out to the public.
\end{cousinweb}

\begin{grammarlens}[title={What you've built}]
One more word for telling, explaining and describing at B1.
\end{grammarlens}

\subsection*{Your turn}
\begin{itemize}
\raggedright
  \item \emph{Say it:} \emph{výpusk}
  \item \emph{Say it:} \emph{výpusk}, once more
  \item \emph{Say it:} a sentence of your own with \emph{výpusk}
  \item \emph{Recall:} say the Russian for to memorize, then the Russian for a piece of news, then say \emph{výpusk} again
\end{itemize}

\begin{tcolorbox}[breakable,colback=teal!4,colframe=teal!35!black,title={Before you move on}]
What does \ru{выпуск} mean? (\textquotedblleft{}A bulletin, an issue, an edition\textquotedblright{}.) Say it once more.
\end{tcolorbox}
\section[\texorpdfstring{\ru{заголовок} (zagolóvok)}{заголовок (zagolóvok)}]{\ru{заголовок} (zagolóvok) — a headline, a title}

\label{lesson:RU-C279-zagolovok}



\begin{quote}
Before the new one: say the Russian for a piece of news, then the Russian for a bulletin.
\end{quote}

\subsection*{You'll want to know: \ru{заголовок}}
\textbf{\ru{заголовок}} (\emph{zagolóvok}) — \textquotedblleft{}a headline, a title\textquotedblright{}.

\begin{quote}
\ru{Заголовок} \ru{в} \ru{газете} \ru{был} \ru{очень} \ru{большой}.
\end{quote}

\begin{quote}
\emph{Zagolóvok v gazéte byl óchen bolshóy.}
\end{quote}

\begin{quote}
The headline in the newspaper was very big.
\end{quote}

Genitive \emph{\ru{заголовка}}.

\begin{cousinweb}
\emph{\ru{заголовок}} is built on \emph{\ru{за}-} and \emph{\ru{голова}} \textquotedblleft{}head\textquotedblright{}, which is inherited from Proto-Slavic \emph{golva} (EtymDB): what stands at the head.
\end{cousinweb}

\begin{grammarlens}[title={What you've built}]
One more word for telling, explaining and describing at B1.
\end{grammarlens}

\subsection*{Your turn}
\begin{itemize}
\raggedright
  \item \emph{Say it:} \emph{zagolóvok}
  \item \emph{Say it:} \emph{zagolóvok}, once more
  \item \emph{Say it:} a sentence of your own with \emph{zagolóvok}
  \item \emph{Recall:} say the Russian for a piece of news, then the Russian for a bulletin, then say \emph{zagolóvok} again
\end{itemize}

\begin{tcolorbox}[breakable,colback=teal!4,colframe=teal!35!black,title={Before you move on}]
What does \ru{заголовок} mean? (\textquotedblleft{}A headline, a title\textquotedblright{}.) Say it once more.
\end{tcolorbox}
\section[\texorpdfstring{\ru{статья} (statyá)}{статья (statyá)}]{\ru{статья} (statyá) — an article}

\label{lesson:RU-C279-statya}



\begin{quote}
Before the new one: say the Russian for a bulletin, then the Russian for a headline.
\end{quote}

\subsection*{You'll want to know: \ru{статья}}
\textbf{\ru{статья}} (\emph{statyá}) — \textquotedblleft{}an article\textquotedblright{}.

\begin{quote}
\ru{Он} \ru{написал} \ru{статью} \ru{о} \ru{нашей} \ru{школе}.
\end{quote}

\begin{quote}
\emph{On napisál statyú o náshey shkóle.}
\end{quote}

\begin{quote}
He wrote an article about our school.
\end{quote}

Genitive \emph{\ru{статьи}}. Also an article of a law: \emph{\ru{статья} \ru{закона}}.

\begin{cousinweb}
\emph{\ru{статья}} is usually traced to the root of \emph{\ru{стать}} \textquotedblleft{}to stand, to become\textquotedblright{}, inherited from Proto-Slavic \emph{stati} (EtymDB), with an old sense \textquotedblleft{}a part in its place\textquotedblright{}; EtymDB lists no parent for \emph{\ru{статья}} itself.
\end{cousinweb}

\begin{grammarlens}[title={What you've built}]
One more word for telling, explaining and describing at B1.
\end{grammarlens}

\subsection*{Your turn}
\begin{itemize}
\raggedright
  \item \emph{Say it:} \emph{statyá}
  \item \emph{Say it:} \emph{statyá}, once more
  \item \emph{Say it:} a sentence of your own with \emph{statyá}
  \item \emph{Recall:} say the Russian for a bulletin, then the Russian for a headline, then say \emph{statyá} again
\end{itemize}

\begin{tcolorbox}[breakable,colback=teal!4,colframe=teal!35!black,title={Before you move on}]
What does \ru{статья} mean? (\textquotedblleft{}An article\textquotedblright{}.) Say it once more.
\end{tcolorbox}
\section[\texorpdfstring{\ru{сенсация} (sensátsiya)}{сенсация (sensátsiya)}]{\ru{сенсация} (sensátsiya) — a sensation}

\label{lesson:RU-C279-sensatsiya}



\begin{quote}
Before the new one: say the Russian for a headline, then the Russian for an article.
\end{quote}

\subsection*{You'll want to know: \ru{сенсация}}
\textbf{\ru{сенсация}} (\emph{sensátsiya}) — \textquotedblleft{}a sensation\textquotedblright{}.

\begin{quote}
\ru{Эта} \ru{новость} \ru{стала} \ru{настоящей} \ru{сенсацией}.
\end{quote}

\begin{quote}
\emph{Éta nóvost stála nastoyáshchey sensátsiyey.}
\end{quote}

\begin{quote}
This news became a real sensation.
\end{quote}

Feminine. The adjective is \emph{\ru{сенсационный}}.

\begin{cousinweb}
\emph{\ru{сенсация}} was borrowed from French \emph{sensation}, built on Latin \emph{sensus} \textquotedblleft{}feeling\textquotedblright{}, from \emph{sentīre} \textquotedblleft{}to feel\textquotedblright{} (Lewis \& Short); EtymDB lists the Russian word with no parent.
\end{cousinweb}

\begin{grammarlens}[title={What you've built}]
One more word for telling, explaining and describing at B1.
\end{grammarlens}

\subsection*{Your turn}
\begin{itemize}
\raggedright
  \item \emph{Say it:} \emph{sensátsiya}
  \item \emph{Say it:} \emph{sensátsiya}, once more
  \item \emph{Say it:} a sentence of your own with \emph{sensátsiya}
  \item \emph{Recall:} say the Russian for a headline, then the Russian for an article, then say \emph{sensátsiya} again
\end{itemize}

\begin{tcolorbox}[breakable,colback=teal!4,colframe=teal!35!black,title={Before you move on}]
What does \ru{сенсация} mean? (\textquotedblleft{}A sensation\textquotedblright{}.) Say it once more.
\end{tcolorbox}
